% Options for packages loaded elsewhere
\PassOptionsToPackage{unicode}{hyperref}
\PassOptionsToPackage{hyphens}{url}
%
\documentclass[
  11pt,
]{article}
\usepackage{amsmath,amssymb}
\usepackage{iftex}
\ifPDFTeX
  \usepackage[T1]{fontenc}
  \usepackage[utf8]{inputenc}
  \usepackage{textcomp} % provide euro and other symbols
\else % if luatex or xetex
  \usepackage{unicode-math} % this also loads fontspec
  \defaultfontfeatures{Scale=MatchLowercase}
  \defaultfontfeatures[\rmfamily]{Ligatures=TeX,Scale=1}
\fi
\usepackage{lmodern}
\ifPDFTeX\else
  % xetex/luatex font selection
\fi
% Use upquote if available, for straight quotes in verbatim environments
\IfFileExists{upquote.sty}{\usepackage{upquote}}{}
\IfFileExists{microtype.sty}{% use microtype if available
  \usepackage[]{microtype}
  \UseMicrotypeSet[protrusion]{basicmath} % disable protrusion for tt fonts
}{}
\makeatletter
\@ifundefined{KOMAClassName}{% if non-KOMA class
  \IfFileExists{parskip.sty}{%
    \usepackage{parskip}
  }{% else
    \setlength{\parindent}{0pt}
    \setlength{\parskip}{6pt plus 2pt minus 1pt}}
}{% if KOMA class
  \KOMAoptions{parskip=half}}
\makeatother
\usepackage{xcolor}
\usepackage[margin=0.8in]{geometry}
\setlength{\emergencystretch}{3em} % prevent overfull lines
\providecommand{\tightlist}{%
  \setlength{\itemsep}{0pt}\setlength{\parskip}{0pt}}
\setcounter{secnumdepth}{-\maxdimen} % remove section numbering
\ifLuaTeX
  \usepackage{selnolig}  % disable illegal ligatures
\fi
\IfFileExists{bookmark.sty}{\usepackage{bookmark}}{\usepackage{hyperref}}
\IfFileExists{xurl.sty}{\usepackage{xurl}}{} % add URL line breaks if available
\urlstyle{same}
\hypersetup{
  pdftitle={Question Conditioned Online Frame Selection for Smart Glasses},
  hidelinks,
  pdfcreator={LaTeX via pandoc}}

\title{Question Conditioned Online Frame Selection for Smart Glasses}
\author{}
\date{}

\begin{document}
\maketitle

IE 598 Final Project Proposal \textbar{} Reinforcement learning under a
frame transmission budget

\section{Objective}

Smart glasses cannot transmit every video frame, and head motion often
blurs the frames they do send. We will learn a policy that decides, frame
by frame, whether to transmit so that a vision--language model (VLM) can
answer a question given before the video starts, using at most $B$
frames. Decisions are causal and irreversible: a skipped frame cannot be
recovered. The server keeps all transmitted frames in temporal order.

\section{Problem formulation}

An episode consists of a question $Q$, a video with $T$ frames,
synchronized IMU measurements, and a reference answer $g$. The clip
length $T$ is not known to the policy in advance. At each time
$t$, the policy chooses $a_t\in\{\mathrm{Transmit},\mathrm{Wait}\}$.
Let $S_t$ be the set of frames transmitted before $t$ ($S_1=\varnothing$)
and $b_t$ the remaining budget:
\[
S_{t+1}=\begin{cases}S_t\cup\{t\},&a_t=\mathrm{Transmit},\\S_t,&a_t=\mathrm{Wait},\end{cases}
\qquad
\sum_{t=1}^{T}\mathbf{1}\{a_t=\mathrm{Transmit}\}\le B .
\]
The task objective is the QA accuracy
$\mathbb{E}[\mathbf{1}\{M(Q,F(S))=g\}]$, where $S=S_{T+1}$, $M$ is a
fixed VLM, and $F(S)$ denotes the selected frames in temporal order.
This objective is too expensive to use as a training signal, so the
MDP we solve instead maximizes a cheap surrogate return (defined
below); QA accuracy is the target this surrogate approximates and
is used for evaluation. Frame count serves as the transmission cost.

\textbf{State and observation.} Transmission decisions do not affect the
recorded video, so we define the full state as
$x_t=(V,Q,g,\{I_j\},t,S_t,b_t)$, where $V$ is the entire video and IMU
stream (which also determines $T$) and $\{I_j\}$ are the evidence
intervals (defined below). The transition is deterministic, the episode
ends when the video ends, and all randomness enters through the initial
state
$x_1\sim P_{x_1}$, the distribution of recorded episodes. Thus $x_t$ is
Markov by construction, and the reward is a function of $(x_t,a_t)$.
The policy, however, observes only a causal function of $x_t$:
\[
y_t=(z_Q,\,z_t,\,q_t,\,m_t,\,h_t,\,b_t,\,t),
\]
where $z_Q$ is a question embedding, $z_t$ features of the current frame
from a lightweight on-device image encoder, $q_t$ its sharpness, $m_t$ a recent IMU summary, and $h_t$ a summary of
past observations and selections; the elapsed time $t$ is observed,
but the remaining time is not. The problem is therefore a POMDP, and
$h_t$ approximates the belief over episodes given the history.

\section{Learning approach and rewards}

The policy is meant to run on the glasses, so every input in $y_t$
comes from cheap computation: a lightweight frozen image encoder for
$z_t$, image sharpness, IMU summaries, and a question embedding $z_Q$
computed once before the video starts. Rewards and the VLM are used
only on the server, for training and evaluation respectively. All
experiments simulate this setting offline on recorded clips, with
features computed once and cached. On these features we train a small
policy network with deep RL (DQN or actor--critic), replaying each
training clip frame by frame so that the policy never sees future
frames. Querying the VLM at the end of every training episode would be
too expensive, so the policy is trained with one of two cheap surrogate
reward functions $r_t(x_t,a_t)$ defined below, maximizing the
undiscounted return $\mathbb{E}\bigl[\sum_{t=1}^{T}r_t\bigr]$. The
fixed VLM $M$ is used only for evaluation on held-out clips. Both
rewards are zero when $a_t=\mathrm{Wait}$; below we give $r_t$ for
$a_t=\mathrm{Transmit}$, with sharpness $q_t\in[0,1]$.

\textbf{Evidence-based reward.} We annotate intervals
$I_1,\ldots,I_m$ that contain the required visual evidence. Let
$\bar q_{t,j}=\max\bigl(\{q_i:i\in S_t\cap I_j\}\cup\{0\}\bigr)$ be the
best sharpness already captured in $I_j$. Then
\[
r_t=\frac{1}{m}\sum_{j:\,t\in I_j}\bigl(q_t-\bar q_{t,j}\bigr)^{+},
\]
so a frame earns reward only if it improves the evidence already
collected. The return equals
$\frac{1}{m}\sum_{j}\max\bigl(\{q_i:i\in S\cap I_j\}\cup\{0\}\bigr)$,
the average best sharpness per evidence interval. Annotations are used
only for rewards and evaluation, never as inputs.

\textbf{CLIP-based reward.} A large frozen image--text encoder
\cite{clip}, used only on the server and distinct from the on-device
encoder, gives relevance
$c_t=\max\bigl(0,\cos(e_t,E_{\mathrm{text}}(Q))\bigr)$ with
$e_t=E_{\mathrm{image}}(f_t)$. To discourage redundant frames, the
reward is discounted by similarity to frames already sent:
\[
r_t=\frac{1}{B}\,q_t c_t\Bigl(1-\max\bigl(\{\cos(e_t,e_i):i\in S_t\}\cup\{0\}\bigr)\Bigr).
\]
This requires no annotation. As with the evidence-based reward, $r_t$
depends on $S_t$, which is part of $x_t$.

\section{Data}

We will record up to 200 first-person clips of 30--60 seconds with
synchronized video and IMU, covering object manipulation, object
relocation, and short action sequences under varied head motion. Each
clip has at least one multiple-choice question. Budgets are 3--5
frames. Data are split 60/20/20 by recording session, keeping all
questions from one recording together. Because the environment is the
unknown episode distribution $P_{x_1}$ and training clips are replayed
many times, performance on training clips overstates generalization;
validation clips are used for model selection and test clips only for
final evaluation.

\section{Evaluation}

Baselines are fixed-interval sampling, causal sharpness and
IMU-assisted thresholds, and a causal question-relevance threshold.
Checkpoints and baseline thresholds are selected by surrogate return on
validation clips. We also ablate the IMU and the question embedding.
For every method, the selected test frames are fed to the fixed VLM
under a common budget and resolution, and we report QA accuracy,
averaged over training seeds. Comparing policies trained with the two
rewards shows which surrogate better serves the downstream task.

\section{Related work and contribution}

RL-based frame selection for video understanding \cite{qin} and
lightweight selectors for frozen VLMs \cite{hornet} already exist, so
question-conditioned selection is not itself our novelty. AdaFrame
\cite{adaframe} learns adaptive frame selection for recognition, and
Koschel et al.\ \cite{koschel} use inertial data for real-time keyframe
selection. Our contribution is a controlled study of irreversible,
budgeted selection on smart glasses: whether question conditioning and
IMU information help, and which surrogate reward better improves QA
accuracy.

\begin{thebibliography}{9}
\bibitem{qin} Y. Qin, H. Li, W. Mu, and Y. He. ``Efficient Frame Selection for
Long Video Understanding via Reinforcement Learning.'' CVPR, pp.
16944--16953, 2026.
\href{https://openaccess.thecvf.com/content/CVPR2026/html/Qin_Efficient_Frame_Selection_for_Long_Video_Understanding_via_Reinforcement_Learning_CVPR_2026_paper.html}{Paper}

\bibitem{hornet} X. Bai, B. Galoaa, and S. Ostadabbas. ``HORNet: Task-Guided
Frame Selection for Video Question Answering with Vision-Language
Models.'' CVPR Workshops, pp. 8755--8764, 2026.
\href{https://openaccess.thecvf.com/content/CVPR2026W/CV4Smalls/html/Bai_HORNet_Task-Guided_Frame_Selection_for_Video_Question_Answering_with_Vision-Language_CVPRW_2026_paper.html}{Paper}

\bibitem{clip} A. Radford et al. ``Learning Transferable Visual Models From
Natural Language Supervision.'' ICML, PMLR 139, pp. 8748--8763, 2021.
\href{https://proceedings.mlr.press/v139/radford21a.html}{Paper}

\bibitem{koschel} A. Koschel, C. Müller, and A. Reiterer. ``Selection of Key
Frames for 3D Reconstruction in Real Time.'' Algorithms, 14(11), 303,
2021. \href{https://www.mdpi.com/1999-4893/14/11/303}{Paper}

\bibitem{adaframe} Z. Wu, C. Xiong, C.-Y. Ma, R. Socher, and L. S. Davis.
``AdaFrame: Adaptive Frame Selection for Fast Video Recognition.'' CVPR,
2019.
\href{https://openaccess.thecvf.com/content_CVPR_2019/html/Wu_AdaFrame_Adaptive_Frame_Selection_for_Fast_Video_Recognition_CVPR_2019_paper.html}{Paper}
\end{thebibliography}
\end{document}
